\chapter{Huber's Further Robust Methods}
\label{sec:p15-huber-remaining-methods}

Huber's examples distinguish sensitivity to tails, to the density
at a median, and to behavior at a rejection boundary. His second
neighborhood construction constrains cumulative probabilities rather
than requiring a contaminated mixture
\cite[Secs.~4, 7 and 9]{Huber1964RobustLocation}.
These are different statistical models. The comparisons below retain
the paper's numerical units and distinguish its asymptotic proposals
from proved results.

\section{Median, trimming and the population being estimated}
\label{sec:p15-huber-median-trimming}

\begin{hbtheorem}{thm:p15-huber-mean-median}{proved=annex:p15.huber.mean.median}{The first two source examples}
For \(n\geq1\) real observations \(x_1,\ldots,x_n\), minimizing
\(\sum_i(x_i-t)^2\) gives their arithmetic mean.
Minimizing \(\sum_i|x_i-t|\) gives the median interval
\[
 [x_{(\lceil n/2\rceil)},x_{(\lfloor n/2\rfloor+1)}],
\]
where \(x_{(j)}\) is the \(j\)th ordered observation.
Its midpoint is one specified median estimator.
For i.i.d.\ observations, the mean is consistent when \(E|X|<\infty\);
if \(EX=0\) and \(EX^2<\infty\), its normalized limit has variance
\(EX^2\). If the distribution function satisfies
\[
 F(0)=1/2,\qquad F(t)=1/2+f_0t+o(|t|),\quad 0<f_0<\infty,
\]
the midpoint sample median is consistent for zero and
\(\sqrt n T_n\) converges to \(N(0,1/(4f_0^2))\).
For the standard normal this variance is \(\pi/2\).
The median claim does not require a finite variance for \(X\).
\end{hbtheorem}

\begin{hbdefinition}{def:p15-huber-hard-trimming}{The source's capped quadratic}
For a fixed residual threshold \(k>0\), put
\[
 \rho^{\rm tr}_k(u)=\tfrac12\min(u^2,k^2),\qquad
 Q^{\rm tr}_n(t)=\sum_i\rho^{\rm tr}_k(x_i-t).
\]
A hard-trimming estimate is a selected global minimizer of this
criterion. It is not defined by discarding a predetermined fraction
of the ordered sample. Its minimizing set can be disconnected;
the midpoint of that set need not be a minimizer.
Away from residuals \(\pm k\), the score is
\(\psi^{\rm tr}_k(u)=u{\bf1}_{\{|u|<k\}}\).
\end{hbdefinition}

\begin{hbtheorem}{thm:p15-huber-trimming-calculation}{proved=annex:p15.huber.trimming.calculation}{Global computation, deletion and a failed consistency inference}
The capped-quadratic criterion has a global minimizer for every
nonempty finite sample. It can be computed by partitioning the line
at all \(x_i\pm k\), minimizing the quadratic on every intervening
interval, and comparing these values with the breakpoint values.
At a minimizer with no residual exactly \(\pm k\), the active set
\(A=\{i:|x_i-t|<k\}\) is nonempty and
\[
                       t=|A|^{-1}\sum_{i\in A}x_i .
\]
Removing any observations outside this active interval retains
that point as a global minimizer. It does not assert that an arbitrary
selection convention remains unchanged.

For \((0,1,2,100)\) and \(k=1\), the capped criterion has precisely
the two minimizers \(1/2,3/2\), each with value \(5/4\);
their midpoint \(1\) has value \(3/2\).
The ordinary mean is \(103/4\), and the midpoint median and
the clipped-linear-score estimate at \(k=1\) are both \(3/2\).

If \(F\) has a density continuous at \(\pm k\), the derivative
at zero of its population trimmed score is
\[
 \lambda'(0)=-D,\qquad
 D=F(k)-F(-k)-k(f(k)+f(-k)).
\]
The formal source variance is
\[
 \frac{\int_{-k}^{k}u^2f(u)\,du}{D^2}.
\]
The derivative identity is proved here. A positive \(D\) alone
does not establish consistency of a global minimizer: for \(k=1\),
the symmetric mixture assigning weights \(1/10,9/20,9/20\)
to uniforms on \([-1/10,1/10]\), \([10-1/10,10+1/10]\),
and \([-10-1/10,-10+1/10]\) has \(D=1/10>0\), but the
expected criterion at \(10\) is smaller than at zero.
\end{hbtheorem}

The source explicitly assumes consistency before its formal trimmed
variance calculation and states that the formula was not proved there.
Its associated conjecture concerns positivity of the denominator's
unsquared expression \cite[p.~79]{Huber1964RobustLocation}.
The preceding counterexample concerns omitting consistency; it is
not a refutation of that conditional conjecture. No present-day
resolution is inferred merely from the 1964 statement.

\begin{hbtheorem}{thm:p15-huber-smooth-stationary-limit}{proved=annex:p15.huber.smooth.stationary.limit}{The separate smooth nonconvex score argument}
Let \(X_i\) be i.i.d., let \(\psi\in C^1(\mathbb R)\) have uniformly
continuous derivative, and assume \(E\psi(X)^2<\infty\),
\(E|\psi'(X)|<\infty\), and \(0<d=E\psi'(X)<\infty\).
Suppose a measurable sequence \(T_n\to0\) in probability satisfies
\(\sum_i\psi(X_i-T_n)=0\).
Then \(E\psi(X)=0\) follows from these assumptions, and
\[
                 \sqrt nT_n\ \Longrightarrow\
                 N(0,E\psi(X)^2/d^2).
\]
This proves the source's smooth stationary-score route with its
integrability made explicit. It does not cover the discontinuous
hard-trimming score or prove consistency of a nonconvex global minimum.
\end{hbtheorem}

\section{The exact median bias comparison}
\label{sec:p15-huber-median-bias}

\begin{hbtheorem}{thm:p15-huber-exact-median-bias}{proved=annex:p15.huber.exact.median.bias}{Contamination bias and indistinguishable locations}
For \(0\leq\varepsilon<1/2\), write \(a=1-\varepsilon\) and
\[
                       c_\varepsilon=\Phi^{-1}(1/(2a)).
\]
Every median of \(a\Phi+\varepsilon H\) belongs to
\([-c_\varepsilon,c_\varepsilon]\); contaminants sufficiently far
to one side attain the corresponding endpoint.
Among all real-valued translation-equivariant location functionals
defined on the contamination model, this is the smallest possible
worst absolute location error. The statement concerns the functional's
population target; a sample estimator needs convergence to that
target before the quantity is its asymptotic bias.
For \(1/2\leq\varepsilon<1\), that worst-error bound is infinite.

The source's flat-central-density construction supplies the smaller
lower bound
\[
 c_{\rm flat}=\frac{\varepsilon}{2(1-\varepsilon)\phi(0)}.
\]
For \(0<\varepsilon<1/2\), \(c_{\rm flat}<c_\varepsilon\);
it is a first-order approximation to the exact median bias, not
its exact upper bound. At \(\varepsilon=0.1,0.01,0.001\), respectively,
the exact median bounds are approximately
\(0.1397103,0.01266008,0.001254569\), in normal standard-deviation units.
\end{hbtheorem}

The source compares clipped-score approximate biases
\(0.17,0.021,0.0027\), for contamination fractions
\(0.1,0.01,0.001\) and respective clipping thresholds
\(1.14,1.95,2.64\). Its comparison with half a nominal standard error,
\(c\leq1/(2\sqrt n)\), leads to reported sample sizes about
\(9,600,35000\), and hence expected contaminant counts about
\(1,6,35\) \cite[pp.~83, 86]{Huber1964RobustLocation}.
These are rounded illustrations of the relation
\(n\leq1/(4c^2)\) for that bias-to-noise ratio. A fixed contamination
bias does not disappear as more observations are taken. The comparison
does not prove that collecting observations beyond those counts is
harmful, nor that such samples identify the contaminating law.

\section{Historical known-scale numerical comparison}
\label{sec:p15-huber-table-one}

Put \(A(k)=2\Phi(k)-1\) and
\[
 B(k)=A(k)-2k\phi(k)+2k^2(1-\Phi(k)).
\]
For \(k>0\), the source's Table I records \(A(k)=E_\Phi\psi_k'\),
\(B(k)=E_\Phi\psi_k^2\), the tuning fraction
\[
 \varepsilon_{\min}(k)=1-
       [A(k)+2\phi(k)/k]^{-1},
\]
and the sharp symmetric-contamination variance bound
\[
 W(k,\varepsilon)
 =\frac{(1-\varepsilon)B(k)+\varepsilon k^2}
        {(1-\varepsilon)^2A(k)^2}.
\]
These formulas and their sharpness are
\cref{thm:p15-huber-normal-variance,thm:p15-huber-least-favorable-normal}.
The \(k=0\) row is their median limit
\(W(0,\varepsilon)=\pi/[2(1-\varepsilon)^2]\); it is not the
literal all-zero score. Its tuning fraction one is likewise a limit.
Every variance in the following three panels is for \(\sqrt nT_n\),
with nominal scale one, and is the source's rounded historical value.
The panels preserve all columns of Table I \cite[pp.~84--85]{Huber1964RobustLocation}.

\begin{center}\small
\begin{longtable}{rrrr}
\caption{Huber Table I: score moments and tuning fraction}
\label{tab:p15-huber-table-one-moments}\\
\(k\)&\(A(k)\)&\(B(k)\)&\(\varepsilon_{\min}\)\\\hline
\endfirsthead
\caption[]{Huber Table I: score moments and tuning fraction (continued)}\\
\(k\)&\(A(k)\)&\(B(k)\)&\(\varepsilon_{\min}\)\\\hline
\endhead
\hline\multicolumn{4}{r}{\emph{continued on next page}}\\
\endfoot
\endlastfoot
0.0&0.0000&0.0000&1.0000\\
0.1&0.0797&0.0095&0.8753\\
0.2&0.1585&0.0358&0.7542\\
0.3&0.2358&0.0758&0.6401\\
0.4&0.3108&0.1265&0.5354\\
0.5&0.3829&0.1851&0.4417\\
0.6&0.4515&0.2491&0.3599\\
0.7&0.5161&0.3160&0.2899\\
0.8&0.5763&0.3840&0.2311\\
0.9&0.6319&0.4511&0.1825\\
1.0&0.6827&0.5161&0.1428\\
1.1&0.7287&0.5777&0.1109\\
1.2&0.7699&0.6352&0.0855\\
1.3&0.8064&0.6880&0.0655\\
1.4&0.8385&0.7358&0.0498\\
1.5&0.8664&0.7785&0.0376\\
1.6&0.8904&0.8160&0.0282\\
1.7&0.9109&0.8487&0.0211\\
1.8&0.9281&0.8767&0.0156\\
1.9&0.9426&0.9006&0.0115\\
2.0&0.9545&0.9205&0.0084\\
2.1&0.9643&0.9371&0.0061\\
2.2&0.9722&0.9507&0.0044\\
2.3&0.9786&0.9617&0.0032\\
2.4&0.9836&0.9705&0.0023\\
2.5&0.9876&0.9776&0.0016\\
2.6&0.9907&0.9831&0.0011\\
2.7&0.9931&0.9873&0.0008\\
2.8&0.9949&0.9906&0.0005\\
2.9&0.9963&0.9931&0.0004\\
3.0&0.9973&0.9950&0.0003\\
\end{longtable}
\begin{longtable}{rrrrrr}
\caption{Huber Table I: smaller contamination fractions}
\label{tab:p15-huber-table-one-small}\\
\(k\backslash\varepsilon\)&0&0.001&0.002&0.005&0.010\\\hline
\endfirsthead
\caption[]{Huber Table I: smaller contamination fractions (continued)}\\
\(k\backslash\varepsilon\)&0&0.001&0.002&0.005&0.010\\\hline
\endhead
\hline\multicolumn{6}{r}{\emph{continued on next page}}\\
\endfoot
\endlastfoot
0.0&1.571&1.574&1.577&1.587&1.603\\
0.1&1.492&1.495&1.498&1.508&1.523\\
0.2&1.423&1.426&1.429&1.438&1.454\\
0.3&1.362&1.365&1.368&1.377&1.393\\
0.4&1.309&1.312&1.315&1.324&1.339\\
0.5&1.263&1.266&1.269&1.278&1.293\\
0.6&1.222&1.225&1.228&1.237&1.252\\
0.7&1.187&1.190&1.193&1.202&1.217\\
0.8&1.156&1.159&1.162&1.172&1.188\\
0.9&1.130&1.133&1.136&1.146&1.162\\
1.0&1.107&1.111&1.114&1.124&1.140\\
1.1&1.088&1.091&1.095&1.105&1.122\\
1.2&1.072&1.075&1.079&1.089&1.107\\
1.3&1.058&1.062&1.065&1.077&1.095\\
1.4&1.047&1.050&1.054&1.066&1.086\\
1.5&1.037&1.041&1.045&1.057&1.078\\
1.6&1.029&1.034&1.038&1.051&1.073\\
1.7&1.023&1.027&1.032&1.046&1.069\\
1.8&1.018&1.023&1.027&1.042&1.066\\
1.9&1.014&1.019&1.024&1.039&1.065\\
2.0&1.010&1.016&1.021&1.038&1.065\\
2.1&1.008&1.014&1.019&1.037&1.066\\
2.2&1.006&1.012&1.018&1.037&1.068\\
2.3&1.004&1.011&1.017&1.037&1.071\\
2.4&1.003&1.010&1.017&1.038&1.074\\
2.5&1.002&1.010&1.017&1.040&1.078\\
2.6&1.002&1.010&1.018&1.042&1.082\\
2.7&1.001&1.010&1.018&1.044&1.087\\
2.8&1.001&1.010&1.019&1.046&1.092\\
2.9&1.001&1.010&1.020&1.048&1.097\\
3.0&1.000&1.011&1.021&1.051&1.103\\
\end{longtable}
\begin{longtable}{rrrrrr}
\caption{Huber Table I: larger contamination fractions}
\label{tab:p15-huber-table-one-large}\\
\(k\backslash\varepsilon\)&0.020&0.050&0.100&0.200&0.500\\\hline
\endfirsthead
\caption[]{Huber Table I: larger contamination fractions (continued)}\\
\(k\backslash\varepsilon\)&0.020&0.050&0.100&0.200&0.500\\\hline
\endhead
\hline\multicolumn{6}{r}{\emph{continued on next page}}\\
\endfoot
\endlastfoot
0.0&1.636&1.741&1.939&2.454&6.283\\
0.1&1.556&1.658&1.853&2.358&6.137\\
0.2&1.485&1.586&1.778&2.276&6.030\\
0.3&1.424&1.524&1.714&2.209&5.962\\
0.4&1.370&1.470&1.659&2.154&5.930\\
0.5&1.324&1.423&1.613&2.111&5.935\\
0.6&1.284&1.384&1.576&2.079&5.976\\
0.7&1.249&1.351&1.546&2.058&6.053\\
0.8&1.220&1.324&1.523&2.047&6.166\\
0.9&1.195&1.302&1.506&2.046&6.317\\
1.0&1.175&1.284&1.495&2.056&6.506\\
1.1&1.158&1.272&1.490&2.072&6.734\\
1.2&1.144&1.263&1.491&2.099&7.003\\
1.3&1.134&1.258&1.497&2.135&7.314\\
1.4&1.126&1.256&1.507&2.180&7.669\\
1.5&1.121&1.258&1.522&2.233&8.069\\
1.6&1.118&1.262&1.542&2.296&8.517\\
1.7&1.116&1.270&1.567&2.367&9.012\\
1.8&1.117&1.280&1.595&2.448&9.558\\
1.9&1.119&1.292&1.628&2.537&10.154\\
2.0&1.122&1.307&1.665&2.635&10.802\\
2.1&1.127&1.324&1.705&2.742&11.501\\
2.2&1.133&1.343&1.750&2.858&12.253\\
2.3&1.140&1.363&1.798&2.982&13.058\\
2.4&1.148&1.386&1.850&3.115&13.914\\
2.5&1.156&1.410&1.905&3.255&14.821\\
2.6&1.166&1.436&1.963&3.405&15.779\\
2.7&1.176&1.463&2.025&3.562&16.787\\
2.8&1.186&1.492&2.090&3.726&17.843\\
2.9&1.198&1.523&2.158&3.899&18.947\\
3.0&1.209&1.554&2.229&4.078&20.098\\
\end{longtable}
\end{center}

\section{A neighborhood defined by cumulative probabilities}
\label{sec:p15-huber-cdf-neighborhood}

\begin{hbdefinition}{def:p15-huber-cdf-neighborhood}{Proper laws and their symmetric vague completion}
For a proper law \(P\), Huber's CDF neighborhood of a normal law is
\[
 \sup_x|P((-\infty,x])-\Phi((x-\xi)/\sigma)|\leq\varepsilon,
 \qquad \sigma>0.
\]
Here \(\xi,\sigma\) are approximate normal parameters; this inequality
does not identify them uniquely. For the symmetric standard problem,
let \(F\) be a symmetric subprobability on the real line, of mass \(m\),
and define its completed cumulative function
\[
 B_F(x)=(1-m)/2+F((-\infty,x]).
\]
It represents equal missing masses at the two infinities.
Define \(\mathcal K_\varepsilon\) by
\(\sup_x|B_F(x)-\Phi(x)|\leq\varepsilon\).
For proper laws this is the original symmetric neighborhood.
Symmetry means \(F(E)=F(-E)\), or
\(B_F(x)+B_F(-x-)=1\); the left limit is necessary at atoms.
\end{hbdefinition}

\begin{hbtheorem}{thm:p15-huber-cdf-domain}{proved=annex:p15.huber.cdf.domain}{Compactness and the dense finite-information class}
For \(0<\varepsilon<1/2\), \(\mathcal K_\varepsilon\) is convex
and vaguely compact, every member has mass at least \(1-2\varepsilon\),
and its finite-information members are vaguely dense.
The construction of those approximants first mixes with the
standard normal to create slack and then smooths; smoothing alone
need not preserve the original CDF bound.
\end{hbtheorem}

\begin{hbtheorem}{thm:p15-huber-cdf-construction}{proved=annex:p15.huber.cdf.construction}{The cosine core, normal middle and exponential tail}
For \(a>0\), let \(z\in(0,\pi/2)\) solve \(2z\tan z=a^2\),
put \(c=2z/a\), and define
\[
 E(a)=\Phi(a)-\tfrac12-
       \frac{\phi(a)}{\cos^2z}
         \left(\frac a2+\frac{\sin(ca)}{2c}\right).
\]
The function \(E\) is continuous and strictly increasing from zero
to \(1/2\). For each \(0<\varepsilon<1/2\), it determines a unique
\(a\). Also a unique \(b>0\) solves
\[
                    \phi(b)/b-(1-\Phi(b))=\varepsilon .
\]
There is exactly one \(\varepsilon_*\in(0,1/2)\) at which \(a=b\).
The following construction applies for \(0<\varepsilon\leq\varepsilon_*\);
the original describes this range approximately as up to \(0.03\).
It does not supply the continuation beyond the crossing.

For \(x\geq0\), reflect the following density and score evenly and
oddly, respectively:
\[
 f_*(x)=
 \begin{cases}
 \phi(a)\cos^2(cx/2)/\cos^2(ca/2),&0\leq x\leq a,\\
 \phi(x),&a<x<b,\\
 \phi(b)e^{-b(x-b)},&x\geq b,
 \end{cases}
\quad
 \psi_*(x)=
 \begin{cases}
 c\tan(cx/2),&0\leq x\leq a,\\
 x,&a<x<b,\\
 b,&x\geq b.
 \end{cases}
\]
The pieces match in value and first derivative for \(f_*\).
The density integrates to one, belongs to \(\mathcal K_\varepsilon\),
and is positive everywhere. Its information is
\[
 I_*=
 \frac{2c^2\phi(a)}{\cos^2(ca/2)}
       \left(\frac a2-\frac{\sin(ca)}{2c}\right)
 +2\{\Phi(b)-\Phi(a)-b\phi(b)+a\phi(a)\}
 +2b\phi(b).
\]
The case \(\varepsilon=0\) is the standard normal limit.
\end{hbtheorem}

\begin{hbtheorem}{thm:p15-huber-cdf-saddlepoint}{proved=annex:p15.huber.cdf.saddlepoint}{The source variational proof and the atomic boundary cases}
For the preceding admissible parameters, \(f_*\) uniquely minimizes
information over \(\mathcal K_\varepsilon\).
Take \(\chi_*\) equal to the ordinary derivative of \(\psi_*\)
away from its joining points, and at each join take the average
of the two one-sided derivatives. If \(a=b\), use the two actual
outer one-sided derivatives at that single join.
For every \(F\in\mathcal K_\varepsilon\),
\[
 A_F=\int\chi_*\,dF,\qquad B_F=\int\psi_*^2\,dF,\qquad
                   2A_F-B_F\geq I_* .
\]
Thus the efficient score gives a saddlepoint in the specified
extended-score domain, including singular \(F\).

For every proper symmetric \(F\) in this class and i.i.d.\ sampling,
the midpoint convex estimator with derivative \(\psi_*\) is
consistent for zero and asymptotically normal with variance
\(B_F/A_F^2\leq1/I_*\). Atoms at the joins are included:
symmetry makes the actual population derivative equal to the
average-derivative integral just defined.
The worst limiting truncated risk is \(1/I_*\), and the
integrated-score lower bound at \(F_*\) makes this minimal among
translation-equivariant estimators in the same ordered risk sense
as \cref{def:p15-huber-symmetric-location-model}.
\end{hbtheorem}

\begin{hbexample}{def:p15-huber-cdf-numbers}{The original Table II and an independently specified construction}
The following are the rounded parameters in Table II
\cite[p.~92]{Huber1964RobustLocation}; they are tuning values,
not sampling uncertainties:
\[
\begin{array}{c|rrr}
\varepsilon&a&b&c\\\hline
0.001&0.65&2.44&1.37\\
0.002&0.75&2.23&1.35\\
0.005&0.91&1.95&1.32\\
0.01&1.06&1.72&1.29\\
0.02&1.24&1.49&1.26\\
0.03&1.34&1.36&1.23
\end{array}
\]
For a construction specified independently of those rounded entries,
take \(z=\pi/6\). Then
\[
 a=\sqrt{\pi/(3\sqrt3)}\simeq0.77756015,\quad
 c=\pi/(3a)\simeq1.34677369,\quad
 \varepsilon=E(a)\simeq0.0023297693 .
\]
The unique tail root is \(b\simeq2.18553979>a\).
Substitution in the proved information formula gives
\(I_*\simeq0.96245777\), hence variance \(1/I_*\simeq1.03900663\).
These decimal values are evaluations of the displayed exact equations;
no error bound for a numerical optimization method is asserted.
\end{hbexample}

\section{Scale through logarithms of squared observations}
\label{sec:p15-huber-log-scale}

\begin{hbdefinition}{def:p15-huber-scale-construction}{Both clipping regimes and the transformed model}
For \(X\ne0\), set \(Y=\log X^2\). Multiplying \(X\) by a positive
\(\sigma\) translates \(Y\) by \(\tau=\log\sigma^2\).
For a standard normal \(X\), the transformed density and its
location score are
\[
 g_Y(u)=\frac{\exp(u/2-e^u/2)}{\sqrt{2\pi}},\qquad
 -g_Y'(u)/g_Y(u)=\tfrac12(e^u-1).
\]
Its cumulative function is \(G_Y(u)=2\Phi(e^{u/2})-1\).
For \(c>0\), let \(\chi_c(u)\) clip this score at \(-c,c\).
The upper tangent is \(u_+=\log(1+2c)\).
If \(c<1/2\), the lower tangent is \(u_-=\log(1-2c)\);
if \(c\geq1/2\), it is \(-\infty\).
The general prototype construction
\cref{thm:p15-huber-general-density} supplies a positive density
and the normalization
\[
 (1-\varepsilon)^{-1}
 =G_Y(u_+)-G_Y(u_-)
       +\{g_Y(u_-)+g_Y(u_+)\}/c,
\]
with the infinite-endpoint term omitted.
The centered transformed contamination model requires \(E_F\chi_c=0\).
It is this specified population target that defines the calibrated
log scale. An arbitrary contaminant does not automatically preserve it.
\end{hbdefinition}

\begin{hbtheorem}{thm:p15-huber-scale-model}{proved=annex:p15.huber.scale.model}{The source's one-sided scale branch and its boundary qualification}
The transformed prototype is smooth and log-concave, so the
preceding normalization and centered minimax result apply for
every \(c>0\). On the branch treated explicitly in Section 10,
\(c\geq1/2\), put \(q=\sqrt{1+2c}\geq\sqrt2\). Then
\[
 (1-\varepsilon)^{-1}
       =A(q)+\frac{2q\phi(q)}{q^2-1}.
\]
For \(F=(1-\varepsilon)\Phi+\varepsilon H\) on the original
\(X\)-scale with \(H\{0\}=0\), centering means
\(E_F\chi_c(\log X^2)=0\) and forces \(H\{|X|\geq q\}=1\).
If \(H\{|X|=q\}=0\), the log-scale estimate from
\(\sum_i\chi_c(\log X_i^2-\log S_n^2)=0\), for i.i.d.\ sampling,
satisfies
\[
 \sqrt n\log(S_n^2/\sigma^2)\Longrightarrow N(0,1/I_Y),\qquad
 I_Y=\frac{1-\varepsilon}{2}\{A(q)-2q\phi(q)\}>0
\]
after a scale transformation by \(\sigma>0\).
Its limiting truncated risk is the minimax value for this centered
log-scale problem. The comparison includes every positive
scale-equivariant estimate allowed to use observation signs,
by \cref{thm:p15-huber-scale-sign-minimax}.
This variance is constant on the source's
strict-tail class \(H\{|X|>q\}=1\).
An atom at \(|X|=q\) still satisfies centering but requires the
one-sided population-slope calculation; it must not be included
silently in that normal-limit assertion.
\end{hbtheorem}

\begin{hbtheorem}{thm:p15-huber-scale-computation}{proved=annex:p15.huber.scale.computation}{Scale equations, zero observations and a worked sample}
For all \(c>0\), put \(\ell=\max(0,1-2c)\) and \(u=1+2c\).
The positive-scale equation, interpreted algebraically also at
zero observations, is
\[
       \sum_{i=1}^n\min\{u,\max(\ell,x_i^2/S^2)\}=n,\qquad S>0.
\]
If \(j\) observations are nonzero, its left side decreases
continuously from \(ju+(n-j)\ell\) as \(S\downarrow0\) to \(n\ell<n\)
as \(S\to\infty\).
If the first limit is below \(n\), no positive solution exists.
If it equals \(n\), a whole interval adjacent to zero consists
of solutions and there is no positive lower root endpoint.
If it exceeds \(n\), the roots form a nonempty bounded closed
interval with positive lower endpoint; when \(c\geq1/2\) there
is exactly one root. When \(c<1/2\), a flat root interval is possible;
the transformed midpoint selection is the geometric mean of
its two scale endpoints.

For \(q=2\) on the one-sided branch the equation is
\(\sum_i\min(x_i^2/S^2,4)=n\).
The sample \((-1,-1,1,1,M)\), \(M>4\), has the unique solution
\(S=2\). For \((0,0,0,1)\), every \(0<S\leq1/2\) solves it.
For \((0,0,0,0,1)\), no positive solution exists.
A declared degenerate scale zero is a separate convention,
not a solution obtained by dividing by zero.
\end{hbtheorem}

\begin{hbtheorem}{thm:p15-huber-scale-gauging}{proved=annex:p15.huber.scale.gauging}{Normal calibration and historical Table III}
For fixed \(q\geq\sqrt2\) and i.i.d.\ standard-normal observations,
the unadjusted scale tends to the unique \(S_\infty\in(0,1)\) satisfying
\[
                  B(qS_\infty)=S_\infty^2.
\]
Put \(q'=qS_\infty\). The normally calibrated estimate
\(S_n'=S_n/S_\infty\) is equivalently defined by
\[
 \frac1n\sum_i\min\{(X_i/S_n')^2,q'^2\}=B(q'),\qquad
                  S_\infty=\sqrt{B(q')},\quad q=q'/\sqrt{B(q')}.
\]
This calibration makes its normal-model limit one. In the earlier
strict-tail contamination model its limit is \(1/S_\infty\);
changing the calibration does not make both targets one.
At \(q=2\), the exact normalization gives
\(\varepsilon\simeq0.02580420\),
\(q'\simeq1.89734505\), and \(S_\infty\simeq0.94867253\).
The worked sample's normally calibrated estimate is
\(2/S_\infty\), whereas its ordinary root mean square
\(\sqrt{(M^2+4)/5}\) grows without bound.
\end{hbtheorem}

Table III gives the following historical rounded calibration values
\cite[p.~95]{Huber1964RobustLocation}. Its last column is the
contamination fraction for which the associated unadjusted \(q\)
is the source's minimax tuning:
\[
\begin{array}{c|rrr}
q'&S_\infty&q=q'/S_\infty&\varepsilon_{\min}\\\hline
1.1&0.760&1.45&0.182\\
1.5&0.882&1.70&0.074\\
2.0&0.960&2.08&0.019\\
2.5&0.989&2.53&0.0038\\
3.0&0.997&3.01&0.0006
\end{array}
\]
Huber's preference for gauging at the normal distribution expresses
a choice of target. His broader expectation of useful behavior
outside the strict-tail class is an interpretation of the method,
not a proved uniform bias guarantee for arbitrary contamination.

\section{Unknown scale: the three original proposals}
\label{sec:p15-huber-joint-proposals}

\begin{hbmethod}{def:p15-huber-three-proposals}{Original equations and their standing}
Write \(\psi(k,u)=\max(-k,\min(u,k))\), \(k>0\).
Proposal 1 fixes \(k\) and seeks
\[
 \frac1n\sum_i\psi(k,(x_i-T)/S)=0,\qquad
 \frac1n\sum_i{\bf1}_{\{|x_i-T|<kS\}}=A(k),\qquad S>0.
\]
The second equation is generally only approximate: its left side
is an integer divided by \(n\). A boundary convention and a target
integer must be specified in using it. For example at \(n=3,k=1\)
its right side is approximately \(0.682689\), not an attainable
exact fraction with denominator three.

Proposal 2 fixes \(k\), sets \(\beta=B(k)\), and solves
\[
 \sum_i\psi(k,(x_i-T)/S)=0,\qquad
 \sum_i\psi(k,(x_i-T)/S)^2=n\beta .
\]
The following finite-sample theorem also allows a prescribed
\(0<\beta<k^2\), but changing \(\beta\) changes the normal calibration.

Proposal 3 chooses \(k\), now in the data's residual units, to minimize
\[
 \frac{\sum_i\psi(k,x_i-T)^2}
      {[\sum_i{\bf1}_{\{|x_i-T|<k\}}]^2}
 \quad\hbox{subject to}\quad\sum_i\psi(k,x_i-T)=0.
\]
The denominator must be positive. The source suggests retaining
more than \(\sqrt n\) interior observations as a safeguard near
\(k=0\). It presents the three procedures as heuristic proposals
and conjectures the adaptive minimax property of Proposal 3.
That conjecture is not an established theorem of this chapter
\cite[pp.~96--97]{Huber1964RobustLocation}.
\end{hbmethod}

\begin{hbtheorem}{thm:p15-huber-joint-finite}{proved=annex:p15.huber.joint.finite}{The exact existence condition, finite computation and iteration}
Suppose \(0<\beta<k^2\) and the largest multiplicity of any observed
value is strictly less than \(n(1-\beta/k^2)\).
Then Proposal 2's two equations have a unique pair \((T,S)\)
with \(S>0\).
For a candidate division into \(m_1\) lower saturated observations,
\(m_2\) interior observations and \(m_3\) upper saturated observations,
put \(\overline x_2=m_2^{-1}\sum_{\rm interior}x_i\).
At the solution \(m_2>0\), and
\[
 \begin{split}
 S^2&=\frac{\sum_{\rm interior}(x_i-\overline x_2)^2}
 {n\beta-\{m_1+m_3+(m_3-m_1)^2/m_2\}k^2},\\
 T&=\overline x_2+kS(m_3-m_1)/m_2.
 \end{split}
\]
At the solution the numerator and denominator are strictly positive.
Enumerating the finitely many consecutive interior blocks in the
ordered sample, discarding empty interiors and nonpositive
denominators, forming each positive-scale candidate, and
retaining exactly those satisfying
\[
 x_i\leq T-kS,\quad |x_i-T|<kS,\quad x_i\geq T+kS
\]
in their respective groups is a finite computation of the unique pair.

The source instead proposes updating the groups from the preceding
\((T,S)\) and repeating these formulas. Provided every update is
defined, this deterministic iteration reaches a fixed point or
eventually repeats periodically; its convergence is not implied
by existence and uniqueness of the equation's solution.
An empty interior set or nonpositive denominator makes that
update invalid. The source explicitly leaves possible cycling
unresolved.

For the generalized calibration \(k=1,\beta=1/2\), the sample
\((-M,-1,0,1,M)\), \(M>2\), has \(T=0,S=2\).
If all observations coincide and \(\beta>0\), no solution exists.
The positivity of \(\beta\) is necessary: at \(\beta=0\), coincident
observations have arbitrary positive scales.
\end{hbtheorem}

\begin{hbtheorem}{thm:p15-huber-joint-regular-limit}{proved=annex:p15.huber.joint.regular.limit}{A full regular-domain justification of the location variance}
Let \(X_i\) be i.i.d.\ from
\(F=(1-\varepsilon)\Phi+\varepsilon H\), with \(0\leq\varepsilon<1\)
and symmetric \(H\),
and let \(0<\beta<k^2\) satisfy
\(\varepsilon<1-\beta/k^2\).
The finite equations have a unique solution eventually almost surely.
There is a unique \(S_0>0\) such that
\[
 E_F\psi(k,X/S_0)^2=\beta,
\]
and \((T_n,S_n)\to(0,S_0)\) almost surely.
Put \(q=kS_0\). If \(F\{|X|=q\}=0\), then
\[
 \sqrt n T_n\Longrightarrow N(0,V),\qquad
 V=\frac{E_F\psi(q,X)^2}{[F((-q,q))]^2}.
\]
No finite moment of \(X\) is required.
The no-boundary-atom premise gives the Gaussian statement here.
The full identified joint law allowing those atoms, and the resulting
limit of the proposed variance estimate, are supplied in
\cref{thm:p15-huber-joint-atomic-limit,thm:p15-huber-joint-law-and-variance}.

For \(n>1\) and a positive denominator, the source's proposed
variance estimate for \(T_n\) is
\[
 \widehat v_n=
 \frac n{n-1}
 \frac{\sum_i\psi(k,(X_i-T_n)/S_n)^2}
      {[\sum_i{\bf1}_{\{|X_i-T_n|<kS_n\}}]^2}S_n^2 .
\]
In this regular domain \(n\widehat v_n\to V\) almost surely.
For finite \(k\), this is not a claim of finite-sample unbiasedness.
When all residuals are unclipped and \(T_n=\overline X\),
the expression reduces to
\(\sum_i(X_i-\overline X)^2/[n(n-1)]\), the usual mean-variance
estimate.
\end{hbtheorem}

\begin{hbtheorem}{thm:p15-huber-joint-variance-bound}{proved=annex:p15.huber.joint.variance.bound}{The calculation behind historical Table IV}
In the preceding regular domain set \(\beta=B(k)\),
\(r=B(k)/k^2\), and assume additionally \(0\leq\varepsilon<r\).
Let \(q_{\max}\) be the unique positive solution
\[
 \frac{B(q_{\max})}{q_{\max}^2}
                  =\frac{r-\varepsilon}{1-\varepsilon}.
\]
Then \(q=kS_0\leq q_{\max}\), and the sharp bound is
\[
 V\leq
 \frac{q_{\max}^2r}{(1-\varepsilon)^2A(q_{\max})^2}
 =W(q_{\max},\varepsilon).
\]
It is attained by symmetric contaminants strictly outside
\([-q_{\max},q_{\max}]\).
The strict inequality \(\varepsilon<r\) is needed for a finite
root in this bound. A nonpositive right side does not define a
finite \(q_{\max}\).
\end{hbtheorem}

These are the historical Table IV values, for the asymptotic variance
of \(\sqrt nT_n\), with normal scale one. The source obtained them
using linear interpolation in Table I
\cite[pp.~99--100]{Huber1964RobustLocation}; the displayed
\(>4\) entries are retained as inequalities, not replaced by invented
decimal values.
For clipping atoms the same sharp upper bound applies to the
variance of the possibly non-Gaussian limiting law by
\cref{thm:p15-huber-joint-atomic-bound}.
\[
\begin{array}{c|rrrrrr}
k\backslash\varepsilon&0&0.010&0.020&0.050&0.100&0.200\\\hline
1.0&1.107&1.138&1.172&1.276&1.490&2.15\\
1.1&1.088&1.120&1.155&1.265&1.495&2.23\\
1.2&1.072&1.105&1.140&1.258&1.506&2.35\\
1.3&1.058&1.093&1.131&1.256&1.526&2.50\\
1.4&1.047&1.084&1.124&1.257&1.557&2.70\\
1.5&1.037&1.077&1.120&1.264&1.592&2.97\\
1.6&1.029&1.072&1.117&1.275&1.639&3.36\\
1.7&1.023&1.068&1.117&1.288&1.697&3.92\\
1.8&1.018&1.066&1.118&1.306&1.769&>4\\
1.9&1.014&1.065&1.121&1.324&1.859&>4\\
2.0&1.010&1.065&1.127&1.353&2.029&>4
\end{array}
\]
The source's exploratory iteration report describes averages around
one or two updates for \(k=1.5\), sample sizes up to 100 and tails
not too heavy, and about two and four updates for Cauchy samples
of sizes 10 and 100. Its variance-correction discussion mentions
exploratory sample sizes 5 and 10. It does not supply replication
counts or uncertainty estimates for these reports; none are inferred
here. These observations do not establish convergence of the iteration.

\section{The \(M_r\) proposal and pair-average medians}
\label{sec:p15-huber-mr}

\begin{hbdefinition}{def:p15-huber-mr-estimator}{The original subset criterion}
Fix \(r\geq1\), \(n\geq r\), and a symmetric real function
\(\rho_r\) of \(r\) residuals. The proposed \(M_r\) criterion is
\[
 Q_{n,r}(t)=\sum_{\substack{I\subseteq\{1,\ldots,n\}\\|I|=r}}
            \rho_r((X_i-t)_{i\in I}).
\]
A selected minimizer is an estimator only when a measurable
selection and an attained minimum have been supplied.
When differentiable in the common shift, put
\(h_t(x_1,\ldots,x_r)=-\partial_t\rho_r(x_1-t,\ldots,x_r-t)\).
Its normalized score is
\[
 U_n(t)=\binom nr^{-1}\sum_{|I|=r}h_t(X_I).
\]
The sum uses distinct indices. Direct enumeration has \(\binom nr\)
terms; the source's computational caution about large \(r\)
refers to this procedure, not a lower bound ruling out faster
algorithms for particular kernels.
\end{hbdefinition}

\begin{hbtheorem}{thm:p15-huber-mr-projection}{proved=annex:p15.huber.mr.projection}{The needed projection and smooth-root theorem}
Let \(X_i\) be i.i.d., \(r\) fixed, \(h_0\) symmetric and measurable,
with \(Eh_0=0\), \(Eh_0^2<\infty\). Define, up to an arbitrary
finite version on an \(F\)-null set,
\[
 p(x)=E[h_0(x,X_2,\ldots,X_r)].
\]
Then
\[
 U_n(0)=\frac r n\sum_i p(X_i)+R_n,\qquad ER_n^2=O(n^{-2}),
\]
and \(\sqrt nU_n(0)\) has limiting variance \(r^2Ep(X)^2\).
Suppose in addition \((t,x)\mapsto h_t(x)\) is jointly measurable,
\(h_t\) is continuously differentiable in \(t\),
\(E(\partial_t h_t|_{t=0})^2<\infty\), and for a deterministic
modulus \(\omega(s)\to0\),
\[
 |\partial_t h_t(x)-\partial_t h_0(x)|\leq\omega(|t|)
\]
for all tuples and sufficiently small \(t\).
Here a derivative with subscript zero denotes evaluation at \(t=0\).
If \(D=-E\partial_t h_0>0\) and a consistent root \(T_n\to0\)
in probability satisfies \(U_n(T_n)=0\), then
\[
 \sqrt nT_n\Longrightarrow
 N\left(0,\frac{r^2Ep(X)^2}{D^2}\right).
\]
Consistency and root existence are premises here.

If \(h_t(x)=h_0(x_1-t,\ldots,x_r-t)\), the coordinate derivatives
are continuous and integrable under \(F^r\), and differentiation
of \(p\) is justified by an integrable dominating bound on the
first coordinate derivative on each compact first-coordinate
interval, then \(D=rEp'(X)\). This recovers the source's formula
\(Ep^2/(Ep')^2\), with the cancellation of \(r\) explicitly justified.
\end{hbtheorem}

\begin{hbtheorem}{thm:p15-huber-pair-median}{proved=annex:p15.huber.pair.median}{The distinct pair-average example and its comparisons}
For \(r=2\), \(\rho_2(u,v)=|u+v|\), minimizing the source criterion
is equivalent to taking a median of
\[
                  Z_{ij}=(X_i+X_j)/2,\qquad i<j.
\]
Take its midpoint median. The diagonal pairs \(i=j\) are not part
of this source version. If the i.i.d.\ observations have a symmetric
density \(f\in L^2(\mathbb R)\), this estimator is consistent
for zero and
\[
 \sqrt nT_n\Longrightarrow
          N\left(0,\frac1{12(\int f^2)^2}\right).
\]
The proof treats this nonsmooth median through its subset empirical
CDF, rather than invoking the smooth-score theorem without its premises.

For symmetric density contamination with \(0\leq\varepsilon<1\),
\(f=(1-\varepsilon)\phi+\varepsilon h\in L^2\),
\[
 V\leq\frac{\pi}{3(1-\varepsilon)^4}.
\]
At \(\varepsilon=0\) equality holds; for \(0<\varepsilon<1\)
each such proper contaminant gives strict inequality, but the
supremum is the displayed bound.
For \(h(x)=\phi(x/3)/3\), the exact comparison is obtained from
\[
 \int f^2=
 \frac{(1-\varepsilon)^2}{2\sqrt\pi}
 +\frac{2\varepsilon(1-\varepsilon)}{\sqrt{20\pi}}
 +\frac{\varepsilon^2}{6\sqrt\pi}.
\]
For the same mixture, Proposal 2 at \(k=1.5\) has its variance
from the regular theorem together with the unrestricted
normal-mixture bridge \cref{thm:p15-huber-zero-atom-mixture-bridge},
where the unique \(q>0\)
solves
\[
 (1-\varepsilon)B(q)+9\varepsilon B(q/3)
                 =q^2B(1.5)/(1.5)^2,
\]
and its denominator is
\(\{(1-\varepsilon)A(q)+\varepsilon A(q/3)\}^2\).
\end{hbtheorem}

Huber's final page reports the following rounded upper bounds
and then the scale-three normal-mixture comparison
\cite[pp.~100--101]{Huber1964RobustLocation}:
\[
\begin{array}{c|rrrrr}
\varepsilon&0&0.001&0.002&0.005&0.010\\\hline
\pi/[3(1-\varepsilon)^4]&1.047&1.051&1.056&1.068&1.090\\
\varepsilon&0.020&0.050&0.100&0.200&0.500\\\hline
\pi/[3(1-\varepsilon)^4]&1.135&1.286&1.596&2.557&16.76
\end{array}
\]
\[
\begin{array}{c|rrrrr}
\varepsilon&0&0.01&0.02&0.05&0.10\\\hline
\text{pair-average median}&1.047&1.07&1.09&1.17&1.31\\
\text{Proposal 2},\ k=1.5&1.037&1.06&1.08&1.16&1.30
\end{array}
\]
All entries concern asymptotic variances of \(\sqrt nT_n\).
The observations underlying the pair averages are independent;
the pair averages themselves are not. Their dependence is why
the projection calculation is needed.
The source associates this example with Hodges and Lehmann's
1963 work and its projection argument with Hoeffding's 1948 work.
The treatment here supplies the needed arguments for Huber's
specified version; it does not claim those separate originals
have thereby received their own full source treatment.
